\documentclass[a4paper,landscape,8pt]{extarticle}
\usepackage[landscape,margin=0.7cm,top=0.5cm,bottom=0.7cm]{geometry}
\usepackage[utf8]{inputenc}
\usepackage[T1]{fontenc}
\usepackage[ngerman]{babel}
\usepackage{helvet}
\renewcommand{\familydefault}{\sfdefault}
\usepackage{xcolor}
\usepackage{tikz}
\usetikzlibrary{positioning}
\usepackage{enumitem}
\usepackage{multicol}
\usepackage{array}
\usepackage{fancyhdr}
\usepackage{amssymb}
\usepackage{eurosym}

% Farben
\definecolor{cloudblue}{HTML}{2196F3}
\definecolor{syncgreen}{HTML}{4CAF50}
\definecolor{storageoran}{HTML}{FF9800}
\definecolor{lightgray}{HTML}{F5F5F5}
\definecolor{bordergray}{HTML}{BBBBBB}

\pagestyle{fancy}
\fancyhf{}
\renewcommand{\headrulewidth}{0pt}
\fancyfoot[C]{\footnotesize Seite \thepage{} von 3}

\setlength{\parindent}{0pt}
\setlength{\parskip}{2pt}

% Kompakte Listen
\setlist{nosep, leftmargin=1em, itemsep=0pt, topsep=1pt}

\begin{document}

% ============================================================================
% SEITE 1: INFORMATIONSBLATT
% ============================================================================

\noindent{\Large\bfseries iCloud: Dein Schließfach im Internet} \hfill
{\footnotesize\textbf{Lernziele:} Cloud verstehen | Speicher prüfen | Sicher fühlen}

\vspace{1mm}
\hrule
\vspace{1mm}

% Drei-Spalten-Layout
\begin{minipage}[t]{0.32\linewidth}
\begin{center}
\colorbox{cloudblue!15}{%
\begin{minipage}{0.95\linewidth}
\vspace{2mm}
\begin{center}
{\Large\bfseries\textcolor{cloudblue}{1. Was ist die Cloud?}}\\[1pt]
{\footnotesize\itshape Speicher im Internet}
\end{center}
\vspace{2mm}
\end{minipage}%
}
\end{center}

\vspace{1mm}

\textbf{Einfach erklärt:}\\
Die Cloud ist wie ein \textbf{Schließfach bei der Bank} -- nicht bei dir zuhause, aber sicher aufbewahrt.

\textbf{Was bedeutet das?}
\begin{itemize}
\item Deine Daten werden auf \textbf{Computern von Apple} gespeichert
\item Du kannst von \textbf{überall} darauf zugreifen
\item Wenn dein Handy kaputt geht: \textbf{Daten sind sicher!}
\end{itemize}

\textbf{iCloud ist Apples Cloud:}
\begin{itemize}
\item Gehört zu deiner \textbf{Apple-ID}
\item \textbf{5 GB kostenlos} dabei
\item Speichert automatisch viele Daten
\end{itemize}

\textbf{Vorteile:}
\begin{itemize}
\item Handy verloren? Daten noch da!
\item Foto auf iPhone, auch auf iPad sichtbar
\item Sicherungskopie automatisch
\end{itemize}

\end{minipage}%
\hfill%
\begin{minipage}[t]{0.32\linewidth}
\begin{center}
\colorbox{syncgreen!15}{%
\begin{minipage}{0.95\linewidth}
\vspace{2mm}
\begin{center}
{\Large\bfseries\textcolor{syncgreen}{2. Was wird gespeichert?}}\\[1pt]
{\footnotesize\itshape Automatisch in der Cloud}
\end{center}
\vspace{2mm}
\end{minipage}%
}
\end{center}

\vspace{1mm}

\textbf{In iCloud werden automatisch gespeichert:}
\begin{itemize}
\item \textbf{Fotos} (wenn aktiviert)
\item \textbf{Kontakte} und Kalender
\item \textbf{Notizen}
\item \textbf{Safari-Lesezeichen}
\item \textbf{Backup} des gesamten Geräts
\end{itemize}

\textbf{So prüfst du deinen Speicher:}
\begin{enumerate}
\item Einstellungen öffnen
\item Oben auf \textbf{deinen Namen} tippen
\item ``iCloud'' antippen
\item Farbiger Balken zeigt Verbrauch
\end{enumerate}

\textbf{Was verbraucht Speicher?}
\begin{itemize}
\item \textcolor{syncgreen}{\textbf{Grün:}} Fotos
\item \textcolor{storageoran}{\textbf{Gelb:}} Backups
\item \textcolor{bordergray}{\textbf{Grau:}} Anderes (Mail, Dokumente)
\end{itemize}

\end{minipage}%
\hfill%
\begin{minipage}[t]{0.32\linewidth}
\begin{center}
\colorbox{storageoran!15}{%
\begin{minipage}{0.95\linewidth}
\vspace{2mm}
\begin{center}
{\Large\bfseries\textcolor{storageoran}{3. Speicher voll?}}\\[1pt]
{\footnotesize\itshape Was tun?}
\end{center}
\vspace{2mm}
\end{minipage}%
}
\end{center}

\vspace{1mm}

\textbf{5 GB sind schnell voll!}\\
Vor allem Fotos und Videos brauchen viel Platz.

\textbf{Optionen:}

\textbf{a) Mehr Speicher kaufen:}
\begin{itemize}
\item \textbf{50 GB} für 0,99 \euro/Monat
\item \textbf{200 GB} für 2,99 \euro/Monat
\item Einstellungen $\rightarrow$ Name $\rightarrow$ iCloud $\rightarrow$ Speicher verwalten
\end{itemize}

\textbf{b) Speicher freigeben:}
\begin{itemize}
\item Alte Backups löschen
\item Nicht alle Apps sichern
\item Fotos aufräumen (siehe 017)
\end{itemize}

\vspace{2mm}
\fcolorbox{cloudblue}{cloudblue!5}{%
\begin{minipage}{0.92\linewidth}
\footnotesize
\textbf{Empfehlung:} 50 GB für 0,99\euro{}/Monat reicht für die meisten. Sicherheit ist den Kaffee wert!
\end{minipage}%
}

\end{minipage}

\vspace{1mm}

% Zusammenfassung
\begin{center}
\fcolorbox{bordergray}{lightgray}{%
\begin{minipage}{0.98\linewidth}
\centering
\vspace{1.5mm}
\textbf{Zusammenfassung:}
\textcolor{cloudblue}{\textbf{Cloud}} = Sicherer Speicher im Internet ~$|$~
\textcolor{syncgreen}{\textbf{iCloud}} = Apples Cloud (5 GB kostenlos) ~$|$~
\textcolor{storageoran}{\textbf{Mehr Speicher}} ab 0,99\euro{}/Monat
\vspace{1.5mm}
\end{minipage}%
}
\end{center}

\newpage

% ============================================================================
% SEITE 2: ÜBUNGEN
% ============================================================================

\begin{center}
{\LARGE\bfseries Übungen: Cloud verstehen}\\[3pt]
{\normalsize Schau dir Seite 1 an und beantworte die Fragen}
\end{center}

\vspace{3mm}
\hrule
\vspace{4mm}

\begin{multicols}{2}

\textbf{\large Teil A: Grundverständnis}

Richtig oder Falsch? Kreuze an.

\vspace{2mm}

\renewcommand{\arraystretch}{1.6}
\begin{tabular}{@{}p{8cm}cc@{}}
& R & F \\
1. Die Cloud ist ein Speicher im Internet. & $\square$ & $\square$ \\
2. iCloud gehört zu meiner Apple-ID. & $\square$ & $\square$ \\
3. iCloud hat unbegrenzten kostenlosen Speicher. & $\square$ & $\square$ \\
4. Wenn mein Handy kaputt geht, sind Cloud-Daten sicher. & $\square$ & $\square$ \\
5. Nur Fotos werden in iCloud gespeichert. & $\square$ & $\square$ \\
\end{tabular}
\renewcommand{\arraystretch}{1}

\vspace{6mm}

\textbf{\large Teil B: Die Schließfach-Analogie}

Verbinde die Begriffe:

\vspace{2mm}

\renewcommand{\arraystretch}{1.8}
\begin{tabular}{@{}p{4cm}p{5cm}@{}}
Cloud & $\circ$ ~~~~ $\circ$ Automatische Kopie ins Schließfach \\
Sync & $\circ$ ~~~~ $\circ$ Schließfach bei der Bank \\
Backup & $\circ$ ~~~~ $\circ$ Komplette Kopie aller Dokumente \\
5 GB & $\circ$ ~~~~ $\circ$ Kleine Aktentasche \\
\end{tabular}
\renewcommand{\arraystretch}{1}

\vspace{6mm}

\textbf{\large Teil C: Was kostet mehr Speicher?}

Fülle die Lücken aus:

\vspace{2mm}

\renewcommand{\arraystretch}{1.8}
\begin{tabular}{@{}p{4cm}p{5cm}@{}}
5 GB kosten: & \dotfill \\
50 GB kosten: & \dotfill pro Monat \\
200 GB kosten: & \dotfill pro Monat \\
\end{tabular}
\renewcommand{\arraystretch}{1}

\columnbreak

\textbf{\large Teil D: Eigenen Speicher prüfen}

Führe diese Schritte auf deinem Gerät durch:

\vspace{3mm}

\textbf{Schritt 1:} Einstellungen öffnen

$\square$ Erledigt

\vspace{3mm}

\textbf{Schritt 2:} Oben auf deinen Namen tippen

$\square$ Erledigt

\vspace{3mm}

\textbf{Schritt 3:} ``iCloud'' antippen

$\square$ Erledigt

\vspace{3mm}

\textbf{Schritt 4:} Speicherbalken anschauen

\vspace{2mm}
Wie viel Speicher nutzt du? \dotfill von 5 GB

\vspace{2mm}
Was verbraucht am meisten? \dotfill

\vspace{5mm}

\textbf{Nachdenken:}\\
Reichen dir 5 GB, oder wäre mehr Speicher sinnvoll?

\vspace{4mm}
\rule{\linewidth}{0.4pt}
\vspace{4mm}
\rule{\linewidth}{0.4pt}

\end{multicols}

\newpage

% ============================================================================
% SEITE 3: CHECKLISTE
% ============================================================================

\begin{center}
{\LARGE\bfseries Checkliste: Das habe ich verstanden}\\[3pt]
{\normalsize Geh die Liste durch und hake ab, was du sicher weißt}
\end{center}

\vspace{3mm}
\hrule
\vspace{6mm}

% Drei Boxen nebeneinander
\begin{minipage}[t]{0.31\linewidth}
\begin{center}
\fcolorbox{cloudblue}{cloudblue!10}{%
\begin{minipage}{0.92\linewidth}
\vspace{3mm}
\begin{center}
{\large\bfseries\textcolor{cloudblue}{Cloud verstehen}}
\end{center}
\vspace{2mm}
\begin{itemize}[label=$\square$, itemsep=5pt]
\item Ich verstehe, was \textbf{Cloud} bedeutet
\item Ich kenne die \textbf{Schließfach-Analogie}
\item Ich weiß, dass \textbf{iCloud} zu Apple gehört
\item Ich verstehe den \textbf{Vorteil}: Daten sind sicher
\end{itemize}
\vspace{3mm}
\end{minipage}%
}
\end{center}
\end{minipage}%
\hfill%
\begin{minipage}[t]{0.31\linewidth}
\begin{center}
\fcolorbox{syncgreen}{syncgreen!10}{%
\begin{minipage}{0.92\linewidth}
\vspace{3mm}
\begin{center}
{\large\bfseries\textcolor{syncgreen}{Was wird gespeichert}}
\end{center}
\vspace{2mm}
\begin{itemize}[label=$\square$, itemsep=5pt]
\item Ich weiß, dass \textbf{Fotos} gespeichert werden
\item Ich weiß, dass \textbf{Kontakte} gespeichert werden
\item Ich kann meinen \textbf{Speicherstand prüfen}
\item Ich verstehe den \textbf{farbigen Balken}
\end{itemize}
\vspace{3mm}
\end{minipage}%
}
\end{center}
\end{minipage}%
\hfill%
\begin{minipage}[t]{0.31\linewidth}
\begin{center}
\fcolorbox{storageoran}{storageoran!10}{%
\begin{minipage}{0.92\linewidth}
\vspace{3mm}
\begin{center}
{\large\bfseries\textcolor{storageoran}{Speicher-Optionen}}
\end{center}
\vspace{2mm}
\begin{itemize}[label=$\square$, itemsep=5pt]
\item Ich weiß: \textbf{5 GB} sind kostenlos
\item Ich kenne die \textbf{Upgrade-Preise}
\item Ich weiß, wie ich \textbf{Speicher spare}
\item Ich kann \textbf{entscheiden}, ob ich mehr brauche
\end{itemize}
\vspace{3mm}
\end{minipage}%
}
\end{center}
\end{minipage}

\vspace{8mm}

% Gesamtverständnis
\begin{center}
\fcolorbox{black}{lightgray}{%
\begin{minipage}{0.95\linewidth}
\vspace{4mm}
\begin{center}
{\large\bfseries Gesamtverständnis}
\end{center}
\vspace{3mm}
\begin{multicols}{2}
\begin{itemize}[label=$\square$, itemsep=6pt]
\item Ich verstehe, was \textbf{Cloud-Speicher} ist
\item Ich weiß, dass meine Daten \textbf{sicher} sind
\item Ich kann meinen \textbf{Speicherstand} prüfen
\item Ich weiß, was \textbf{mehr Speicher} kostet
\item Ich fühle mich \textbf{sicherer} mit iCloud
\end{itemize}
\end{multicols}
\vspace{3mm}
\end{minipage}%
}
\end{center}

\vspace{8mm}

% Notfall-Notizen
\begin{center}
\fcolorbox{bordergray}{white}{%
\begin{minipage}{0.95\linewidth}
\vspace{4mm}
\begin{center}
{\large\bfseries Meine iCloud-Notizen}\\[3pt]
{\footnotesize (Fülle das aus, um den Überblick zu behalten)}
\end{center}
\vspace{4mm}
\renewcommand{\arraystretch}{2}
\begin{tabular}{@{}p{0.45\linewidth}p{0.45\linewidth}@{}}
\textbf{Mein aktueller Speicher:} & \textbf{Meine Entscheidung:} \\
Genutzt: \dotfill von 5 GB & $\square$ 5 GB reichen mir \\[2mm]
Am meisten: \dotfill & $\square$ Ich möchte mehr (50 GB) \\[2mm]
& $\square$ Ich muss noch nachdenken \\
\end{tabular}
\renewcommand{\arraystretch}{1}
\vspace{4mm}
\end{minipage}%
}
\end{center}

\vfill

\begin{center}
\footnotesize\textit{Stand: \today{} -- Cloud-Speicher schützt deine Daten. Bei Handy-Verlust: Alles noch da!}
\end{center}

\end{document}
